\documentclass[11pt]{article}
\usepackage[a4paper, top=18mm, bottom=18mm, left=20mm, right=20mm]{geometry}
\usepackage[T2A]{fontenc}
\usepackage[utf8]{inputenc}
\usepackage[ukrainian]{babel}
\usepackage{graphicx}
\setlength{\parindent}{0pt}
\setlength{\parskip}{4pt}
\pagestyle{empty}
\begin{document}
%begin
{\large\textbf{Контрольна робота}}

\textbf{Варіант \No\,1}\hfill \textit{ПІБ:}\,\underline{\hspace{60mm}}
\vspace{4mm}

%beginitem
1. Що буде виведено за виконання фрагменту коду?

%code
try:
    k = 2
    while k<=8:
        k += 2
    print(k, end=';')
except:
    print('Z')

%code

%enditem
%beginitem
2. %goodpath:Scripts/2018/Mod2018_1/03-good.txt
Відмітити все, що є коректним оператором С++:
( 1 ) cout>>3;

( 2 ) 'a'=3;

( 3 ) int f(int n);

( 4 ) x<y<z

%enditem
%beginitem
3. Що буде виведено за виконання фрагменту коду?
Записати ВИРАЗ, значенням якого є множина, що складається з кубівцілих чисел від 13 до 2003 (включно), що не діляться на жодне з чисел 2, 3.
%enditem
%beginitem
4. Що буде виведено за виконання фрагменту коду?

a,b,c=2,3,8
def h(b):
global c    a=4
    b=1
    c=5
    return a+b+c

a,b,c=5,4,9
print(h(a),a,b,c)

%enditem
%beginitem
5. Що буде виведено за виконання фрагменту коду?

\begin{center}
	\adjustimage{max size={0.8\linewidth}{0.8\paperheight}}
	{../BinTree/trees_exp/162.png}
\end{center}
Указати послідовність міток вершин дерева, зображеного на рис., що утвориться в результаті обходу дерева в оберненому порядку. Мітки записувати через пробіл.\par Обхід дерева починається з кореня (на рис. це верхня вершина).
%answer:162 оберненому

%enditem
%beginitem
6. Що буде виведено за виконання фрагменту коду?
Записати ВИРАЗ, значенням якого є список, що складається з кубівцілих чисел від 13 до 2003 (включно), що не діляться на жодне з чисел 2, 3, записаних в обернено\-му порядку.

%enditem
%beginitem
7. Що буде виведено за виконання фрагменту коду?
def h(a,b):
      c=80
      if a>-2: c=2
      elif b>=-5: return 3
      else: return 8
      return c
   print(h(-6,-6))
%enditem
%beginitem
8. Що буде виведено за виконання фрагменту коду?

%code
c=['0','1','2','3']
del c[-7:-7]
print(c)
%code


%enditem
%beginitem
9. Що буде виведено за виконання фрагменту коду?
Задано тип вузла списку
struct Node \{int n; Node *prev, *next;\};
У списку зберігаються послідовні цілі числа від -49 до 36. Вказівник p встановлено на вузол, в якому зберігається число -6.
Чому дорівнює значення p->next->next->next->next->next->n ?
%enditem
%beginitem
10. Що буде виведено за виконання фрагменту коду?

%code
print(4 >= 4.0 < 5.0 is False)
%code

%enditem
%beginitem
11. Що буде виведено за виконання фрагменту коду?

print('R', end=' ')
k=8
while k>=2:
    k+=-2
print(k)


%enditem
%beginitem
12. Що буде виведено за виконання фрагменту коду?

a,b,c=8,9,7
def h(a,b=6,c):
    print(a,b,c,end="")

h(2,c=3)
print(a,b,c)


%enditem
%beginitem
13. Що буде виведено за виконання фрагменту коду?

%code
a,b,c=2,3,8
def h(b):
global c    a=4
    b=1
    c=5
    return a+b+c

a,b,c=5,4,9
print(h(a),a,b,c)
%code

%enditem
%beginitem
14. Що буде виведено за виконання фрагменту коду?
Написати програму, що запитує у користувача значення дійсних параметрів $x$ та $y$, обчислює для них значення виразу 
$$\frac{\sqrt x - 99}{(\ctg x - 0.4) (y - 82)}$$ 
та повiдомляє введенi данi та результат обчислення.


%enditem
%beginitem
15. Що буде виведено за виконання фрагменту коду?
a = 2
b = 3
c = 8

def h(b):
global c    a *= 4
    b = 1
    c = 5
    return a + b + c

a = 5
b = 4
c = 9

try:
    print(h(a), a, b, c, sep=';')
except:
    print('Z', a, b, c, sep=';')

%enditem
%beginitem
16. Що буде виведено за виконання фрагменту коду?
%code

def f(n):
    print("f", end="");
    return n>-1

print(f(-6) or f(-9))
%code

%enditem
%beginitem
17. Що буде виведено за виконання фрагменту коду?
try:
    k = 2
    while k<=8:
        k += 2
    print(k, end=';')
except:
    print('Z')

%enditem
%beginitem
18. Що буде виведено за виконання фрагменту коду?
False>"4j"
%enditem
%beginitem
19. Що буде виведено за виконання фрагменту коду?
not4>4.0 or not5.0>=False or 5<=3
%enditem
%beginitem
20. Що буде виведено за виконання фрагменту коду?

%code
c=['0','1','2','3']
c.insert(-3, '12')
print(c)
%code


%enditem
%beginitem
21. Що буде виведено за виконання фрагменту коду?

print('R', end=' ')
c=('0','1','2','3','4','5','6','7','8','9')
print(c[-12:-12:2])


%enditem
%beginitem
22. Що буде виведено за виконання фрагменту коду?

%code
#include <iostream>

int h(int a, int b){
    int c = 80;
    if (a > -2)
        c = 2;
    else if (b >= -5)
         return 3;
    else 
        return 8;
    return c;
}

int main(){
    std::cout << h(-6, -6);
    return 0;
}
%code

%enditem
%beginitem
23. Що буде виведено за виконання фрагменту коду?
%code
for c in range(-12,-12,2):
    if c>5:
        break
        print(c,end=' ')
    if c>5:
        break
else:
    print(c, end=' ')
print(c, end=' ')
%code

%enditem
%beginitem
24. Що буде виведено за виконання фрагменту коду?
Записати ВИРАЗ, значенням якого є список, що складається з цілих чисел від 13 до 2003 (включно), причому спочатку за спаданням записано всі, що діляться на 5, а потім решту за спаданням.
%enditem
%beginitem
25. Що буде виведено за виконання фрагменту коду?

\begin{center}
	\adjustimage{max size={0.8\linewidth}{0.8\paperheight}}
	{../BinTree/trees_exp/162.png}
\end{center}
Указати послідовність міток вершин дерева, зображеного на рис., що утвориться в результаті обходу дерева в оберненому порядку. Мітки записувати через пробіл.\par Алгоритм починає роботу з кореня (на рис. це верхня вершина).
%answer:162 оберненому

%enditem
%beginitem
26. Що буде виведено за виконання фрагменту коду?
%code
#include <iostream>
using namespace std;

int f(int &x, int &y){
    x = 4;
    y-= 1;
    return y;
}

int main(){
    int a = 7, b = 3;
    b = f(b, b);
    cout << a << ":" << b <<':';
    {
        int a = 2, b = 6;
        cout << ((b>6) || ((a-=2) > 6)) << ':';
        cout << a << ":" << b << ':';
    }
    {
        inta = 9;
        cout << a << ':';
    }
    cout << a << ":" << b << endl;
    return 0;
}
%code


%enditem
%beginitem
27. Що буде виведено за виконання фрагменту коду?
try:
    for c in range(-12, -12, 2):
        if c>5:
            break
        print(c, end=';')
        if c>5:
            break
    else:
        print(c,end=';')
    print(c)
except:
    print('ERR')

%enditem
%beginitem
28. Що буде виведено за виконання фрагменту коду?
def f(n):
    print('f', end='')
    return n>-1

try:
    print(f(-6) or f(-9))
except:
    print('ERR')

%enditem
%beginitem
29. Що буде виведено за виконання фрагменту коду?

%code
print(6//6//6*6)
%code

%enditem
%beginitem
30. Що буде виведено за виконання фрагменту коду?

%code
cout << (!4>4.0 && !5.0>= false && 5<=3);
%code

%enditem
%beginitem
31. Що буде виведено за виконання фрагменту коду?
c=('a','b','c','d','e',0,1,2,4); print(c[-12])
%enditem
%beginitem
32. Що буде виведено за виконання фрагменту коду?
Задано тип вузла списку struct Node {int n; Node *prev, *next;};
У списку зберігаються послідовні цілі числа від -49 до 36. Вказівник p встановлено на вузол, в якому зберігається число -6.
Чому дорівнює значення p->next->next->next->next->next->n ?
%enditem
%beginitem
33. Що буде виведено за виконання фрагменту коду?

%code
a = 2
b = 3
c = 8

def h(b):
global c    a=4
    b=1
    c=5
    return a+b+c

a, b, c = 5,4,9
try:
    print(h(a), a, b, c, sep=';')
except:
    print('Z', a, b, c, sep=';')
%code

%enditem
%beginitem
34. Що буде виведено за виконання фрагменту коду?

print('R', end=' ')
k=2
while k<=8:
    k+=2
print(k)


%enditem
%beginitem
35. Що буде виведено за виконання фрагменту коду?

%code
try:
    k = 8
    while k>=2:
        k += -2
    print(k, end=';')
except:
    print('Z')

%code

%enditem
%beginitem
36. Що буде виведено за виконання фрагменту коду?

%code
cout << (6 % 6 % 6 % 6);
%code

%enditem
%beginitem
37. Що буде виведено за виконання фрагменту коду?

print('R', end=' ')
k=8
while k>=2:
    k+=-2
print(k)


%enditem
%beginitem
38. Що буде виведено за виконання фрагменту коду?

print('R', end=' ')
c=['0','1','2','3','4']
del c[-7:-7]
print(c)


%enditem
%beginitem
39. Що буде виведено за виконання фрагменту коду?
%code
if (2 > 2)
    cout << "v";
else
    cout << "y";
%code


%enditem
%beginitem
40. Що буде виведено за виконання фрагменту коду?
Рядки журналу через двокрапку містять ім'я користувача (ідентифікатор), час події,тип події, код події, наприклад
\begin{verbatim}
s="username : 2022-05-05 13:26 : ERROR : 404"
\end{verbatim}
у коді 
\begin{verbatim}
res = re.fullmatch('???', s) 
s = ''
code_str = ??? 
\end{verbatim}
чим треба замінити кожен з ???, щоб значенням code\_str був код події? 



%enditem
%beginitem
41. Що буде виведено за виконання фрагменту коду?

%code
a = 2
b = 3
c = 8

def h(b):
global c    a*=4
    b=1
    c=5
    return a+b+c

a, b, c = 5, 4, 9
try:
    print(h(a), a, b, c, sep=';')
except:
    print('Z', a, b, c, sep=';')
%code

%enditem
%beginitem
42. Що буде виведено за виконання фрагменту коду?
try:
    res = False>"4j"
    if isinstance(res, bool):
        print(f'bool;{res}')
    elif isinstance(res, int):
        print(f'int;{res}')
    elif isinstance(res, float):
        print(f'float;{res:.2f}')
    else:
        print('???')
except:
    print('ERR')

%enditem
%beginitem
43. %goodpath:Scripts/2019/Mod2019_1/Lex/03-good.txt
Відмітити все, що є коректним оператором С++:
( 1 ) n=2; n--

( 2 ) (2+2j)//2

( 3 ) x=1//3

( 4 ) x<3 or y<7

%enditem
%beginitem
44. Що буде виведено за виконання фрагменту коду?
Знайти кількість відношень на n-елементній множині, які неантирефлексивні або неантисиметричні
%enditem
%beginitem
45. Що буде виведено за виконання фрагменту коду?
Скільки 2017-елементних підмножин множини натуральних від 1 до 20172017 містять принаймні одне число, що менше 172002
%enditem
%beginitem
46. Що буде виведено за виконання фрагменту коду?
try:
    print(2, end=';')
    print(int('c3'), end=';')
    print(8, end=';')
except ZeroDivisionError:
    print(5, end=';')
except Exception:
    print(4, end=';')
except:
    print('ERR', end=';')
else:
    print(9, end=';')
finally:
    print(7)

%enditem
%beginitem
47. Що буде виведено за виконання фрагменту коду?

%code
print(-2**-4-1)
%code

%enditem
%beginitem
48. Що буде виведено за виконання фрагменту коду?
В умовах попередньої задачі було виконано p->next=p->prev->prev;

Чому дорівнює значення p->next->next->next->next->n ?
%enditem
%beginitem
49. Що буде виведено за виконання фрагменту коду?

print('R', end=' ')
k=2
while k<=8:
    k+=2
print(k)


%enditem
%beginitem
50. Що буде виведено за виконання фрагменту коду?

def f(n):
    print("f", end="");
    return n>-1

print(f(-6) or f(-9))

%enditem
%beginitem
51. Що буде виведено за виконання фрагменту коду?
Записати ВИРАЗ, значенням якого є множина, що складається з цілих чисел від 13 до 2003 (включно), що діляться на 5.
%enditem
%beginitem
52. Що буде виведено за виконання фрагменту коду?

%code
print('R', end=' ')
c=['0','1','2','3','4']
del c[-7:-7]
print(c)
%code


%enditem
%beginitem
53. Що буде виведено за виконання фрагменту коду?
Рядки журналу через = містять ім'я користувача (ідентифікатор), час події,тип події, код події, наприклад
\begin{verbatim}
s="username = 2022-05-05 13:26 = ERROR = 404"
\end{verbatim}
у коді 
\begin{verbatim}
res = re.fullmatch('???', s) 
s = ''
code_str = ??? 
\end{verbatim}
чим треба замінити кожен з ???, щоб значенням code\_str був код події? 



%enditem
%beginitem
54. Що буде виведено за виконання фрагменту коду?

%code
print(False>"4j")
%code

%enditem
%beginitem
55. Що буде виведено за виконання фрагменту коду?

%code
c = ('a','b','c','d','e',0,1,2,4)
print(c[-12])
%code


%enditem
%beginitem
56. Що буде виведено за виконання фрагменту коду?

%code
#include <iostream>
using namespace std;

int a = 2, b = 3, c = 8;

int h(int &b){
int c;    a *= 4;
    b = 1;
    c = 5;
    return a + b + c;
}

int main(){
    inta = 5;
    intb = 4;
    intc = 9;
    cout << h(a) << ':';
    cout << a << ':' << b << ':' << c;
    return 0; 
}
%code

%enditem
%beginitem
57. Що буде виведено за виконання фрагменту коду?

print('R', end=' ')
c=['0','1','2','3','4']
c.insert(-4, '12')
print(c)


%enditem
%beginitem
58. %goodpath:Scripts/2019/Mod2019_1/Lex/01-good.txt
Відмітити все, що є однією правильною лексемою:
( 1 ) x*y

( 2 ) -56

( 3 ) "a'a"

( 4 ) **

%enditem
%beginitem
59. Що буде виведено за виконання фрагменту коду?
Змінні \kw{next} та \kw{prev} вузла списку позначають наступний та попередній за ним вузол відповідно. Починаючи з голови, у список записано послідовні цілі числа від~$-49$ до~$36$. Нехай змінна \kw{p} позначає вузол, що зберігає значення~$-6$.
Яке число зберігається у вузлі \kw{p.next.next.next.next.next} ?

%enditem
%beginitem
60. Що буде виведено за виконання фрагменту коду?

%code
print(4>=4.0>=5.0>=False)
%code

%enditem
%beginitem
61. Що буде виведено за виконання фрагменту коду?

print('R', end=' ')
c=('0','1','2','3','4','5','6','7','8','9')
print(c[-12:-12:2])


%enditem
%beginitem
62. Що буде виведено за виконання фрагменту коду?

%code
print('R', end=' ')
for c in range(-7,-7,2):
    if c>5:
        break
        print(c, end=' ')
    if c>5:
        break
    else:
        print(c, end=' ')
    print(c, end=' ')
%code

%enditem
%beginitem
63. %goodpath:Scripts/2018/Mod2018_1/01-good.txt
Відмітити все, що є однією правильною лексемою:
( 1 ) <>

( 2 ) x*y

( 3 ) FIFO

( 4 ) ||

%enditem
%beginitem
64. Що буде виведено за виконання фрагменту коду?
Записати ВИРАЗ, значенням якого є множина, що складається з цілих чисел від 13 до 2003 (включно), що діляться на 5.
%enditem
%beginitem
65. Що буде виведено за виконання фрагменту коду?
for c in range(-12,-12,2):
    if c>5:
        break
        print(c,end=' ')
    if c>5:
        break
else:
    print(c, end=' ')
print(c, end=' ')

%enditem
%beginitem
66. Що буде виведено за виконання фрагменту коду?

a,b,c=2,3,8
def h(b):
global c    a*=4
    b=1
    c=5
    return a+b+c

a,b,c=5,4,9
print(h(a),a,b,c)

%enditem
%beginitem
67. Що буде виведено за виконання фрагменту коду?
try:
    res = 6//6//6*6
    if isinstance(res, bool):
        print(f'bool;{res}')
    elif isinstance(res, int):
        print(f'int;{res}')
    elif isinstance(res, float):
        print(f'float;{res:.2f}')
    else:
        print('???')
except:
    print('Z')

%enditem
%beginitem
68. Що буде виведено за виконання фрагменту коду?
Записати ВИРАЗ, значенням якого є словник, ключами якого є цілі числа від 13 до 2003 (включно), що не діляться на жодне з чисел 2, 3, а ключам відповідають їх квадратні корені 

%enditem
%beginitem
69. Що буде виведено за виконання фрагменту коду?
Записати ВИРАЗ, значенням якого є список, що складається з кубівцілих чисел від 13 до 2003 (включно), що не діляться на жодне з чисел 2, 3, записаних в оберненому порядку.
%enditem
%beginitem
70. Що буде виведено за виконання фрагменту коду?

%code
for c in range(-7,-7+5,2):
    if c>5:
        break
    print(c, end=' ')
    c=-7
    if c>5:
        break
else:
    print(c, end=' ')
print(c, end=' ')
%code

%enditem
%beginitem
71. Що буде виведено за виконання фрагменту коду?
6//6//6*6
%enditem
%beginitem
72. Що буде виведено за виконання фрагменту коду?
try:
    res = -2**-4-1
    if isinstance(res, bool):
        print(f'bool;{res}')
    elif isinstance(res, int):
        print(f'int;{res}')
    elif isinstance(res, float):
        print(f'float;{res:.2f}')
    else:
        print('???')
except:
    print('ERR')

%enditem
%beginitem
73. Що буде виведено за виконання фрагменту коду?

%code
print('R', end=' ')
c = ('a','b','c','d','e',0,1,2,4)
print(c[-12])
%code


%enditem
%beginitem
74. Що буде виведено за виконання фрагменту коду?

%code
False>"4j"
%code

%enditem
%beginitem
75. Що буде виведено за виконання фрагменту коду?

%code
#include <iostream>
using namespace std;

bool f(int n){
    cout<<"f";
    return n>-1;
}

int main(){
    cout<<(f(-6) && f(-9));
    return 0;
}
%code

%enditem
%beginitem
76. Що буде виведено за виконання фрагменту коду?

%code
#include <iostream>

int h(int c){
    int w = 80;
    if (c > -2) 
        w = 2;
    else if (c >= -5)
         return 3;
    else
         return 8;
    return w;
}

int main(){
    std::cout << h(-6);
    return 0;
}
%code


%enditem
%beginitem
77. Що буде виведено за виконання фрагменту коду?

\begin{center}
	\adjustimage{max size={0.8\linewidth}{0.8\paperheight}}
	{../BinTree/trees_exp/70.png}
\end{center}
Указати послідовність міток вершин дерева, зображеного на рис., що утвориться в результаті обходу дерева в оберненому порядку. ̳тки записувати через пробіл.\par Обхід дерева починається з кореня (на рис. це верхня вершина).
%answer:70 оберненому
%enditem
%beginitem
78. Що буде виведено за виконання фрагменту коду?

print('R', end=' ')
for c in range(-12,-12,2):
    if c>5:
        break
        print(c, end=' ')
    if c>5:
        break
    else:
        print(c, end=' ')
    print(c)

%enditem
%beginitem
79. Що буде виведено за виконання фрагменту коду?
Записати ВИРАЗ, значенням якого є множина, що складається з кубівцілих чисел від 13 до 2003 (включно), що не діляться на жодне з чисел 2, 3.

%enditem
%beginitem
80. Що буде виведено за виконання фрагменту коду?
a,b,c=2,3,8
   def h(b):
     global c     a=4
     b=1
     c=5
     return a+b+c
   a,b,c=5,4,9
   print(h(a),a,b,c)
%enditem
%beginitem
81. Що буде виведено за виконання фрагменту коду?
for c in range(-12,-12,2):
    if c>5:
        break
        print(c,end=' ')
    if c>5:
        break
    else:
        print(c,end=' ')
    print(c)

%enditem
%beginitem
82. Що буде виведено за виконання фрагменту коду?
%Оригинал был утерян. Требует отладки!
%code
__c = 0

class C:
    __c = 1
    
    def __init__(self):
        self.__c = 2
        __c = 3
        C.__c = 4

obj = C()
C.__c = 5
try:
    print(__c, end=' ')
    print(obj.__c, end=' ')
    print(C.__c, end=' ')
    print(obj.__c, end=' ')
    print(obj._C__c)
except:
    print('ERR')
%code

%enditem
%beginitem
83. Що буде виведено за виконання фрагменту коду?

\begin{center}
	\adjustimage{max size={0.8\linewidth}{0.8\paperheight}}
	{../15BinTree/trees_exp/162.png}
\end{center}
Указати послідовність міток вершин дерева, зображеного на рис., що утвориться в результаті обходу дерева в оберненому порядку. Мітки записувати через пробіл.\par Алгоритм починає роботу з кореня (на рис. це верхня вершина).
%answer:162 оберненому

%enditem
%beginitem
84. Що буде виведено за виконання фрагменту коду?
$a$ є цілочисловим масивом типу $numpy.ndarray$ розмірів $4{\times}4$. Сума елементів масиву $a$ дорівнює 23. Чому дорівнює сума елементів масиву $a$ після виконання 
\begin{verbatim}
a -= 23
\end{verbatim}


Якщо код некоректний, то в якості відповіді зазначити ERROR.



%enditem
%beginitem
85. Що буде виведено за виконання фрагменту коду?
c=['0','1','2','3','4']; del c[-7:-7]; print(c)
%enditem
%beginitem
86. Що буде виведено за виконання фрагменту коду?
Рядки журналу через = містять ім'я користувача (ідентифікатор), час події,тип події, код події, наприклад
\begin{verbatim}
s="username = 2022-05-05 13:26 = ERROR = 404"
\end{verbatim}
у коді 
\begin{verbatim}
res = re.fullmatch('???', s) 
s = ''
code_str = ??? 
\end{verbatim}
чим треба замінити кожен з ???, щоб значенням code\_str був код події? 



%enditem
%beginitem
87. Що буде виведено за виконання фрагменту коду?
-2**-4-1
%enditem
%beginitem
88. Що буде виведено за виконання фрагменту коду?

%code
#include <iostream>
using namespace std;

int h(int a, int b){
    int c = 80;
    if (a > -2)
        c = 2;
    else if (b >= -5)
         return 3;
    else 
        return 8;
    return c;
}

int main(){
    cout << h(-6, -6);
    return 0;
}
%code

%enditem
%beginitem
89. Що буде виведено за виконання фрагменту коду?
Написати програму, що запитує у користувача значення дійсних параметрів $x$ та $y$, обчислює для них значення виразу та повiдомляє введенi данi та результат обчислення.
$$\frac{\sqrt x - 99}{(\ctg x - 0.4) (y - 82)}$$ 


%enditem
%beginitem
90. Що буде виведено за виконання фрагменту коду?
Записати ВИРАЗ, значенням якого є список, що складається з цілих чисел від 13 до 2003 (включно), причому спочатку за спаданням записано всі, що діляться на 5, а потім решту за спаданням.

%enditem
%beginitem
91. Що буде виведено за виконання фрагменту коду?

%code
a,b,c=2,3,8
def h(b):
    a*=4
    b=1
    c=5
    return a+b+c

a,b,c=5,4,9
print(h(a),a,b,c)
%code

%enditem
%beginitem
92. Що буде виведено за виконання фрагменту коду?
False>"4j"
%enditem
%beginitem
93. Що буде виведено за виконання фрагменту коду?
c=['0','1','2','3','4']; del c[-7:-7]; print(c)
%enditem
%beginitem
94. Що буде виведено за виконання фрагменту коду?
def h(c):
      w=80
      if c>-2: w=2
      elif c>=-5: return 3
      else: return 8
      return w
   print(h(-6))
%enditem
%beginitem
95. Що буде виведено за виконання фрагменту коду?
try:
    k = 8
    while k>=2:
        k += -2
    print(k, end=';')
except:
    print('Z')

%enditem
%beginitem
96. Що буде виведено за виконання фрагменту коду?

%code
cout << (4 > 4.0 >= 5.0 <= false);
%code

%enditem
%beginitem
97. Що буде виведено за виконання фрагменту коду?
k=8
   while k>=2: k+=-2
   print(k)
%enditem
%beginitem
98. Що буде виведено за виконання фрагменту коду?

%code
cout << 6 % 6 % 6 % 6;
%code

%enditem
%beginitem
99. Що буде виведено за виконання фрагменту коду?
c=['0','1','2','3','4']; c.insert(-4,'12'); print(c)
%enditem
%beginitem
100. Що буде виведено за виконання фрагменту коду?
Записати ВИРАЗ, значенням якого є список, що складається з цілих чисел від 13 до 2003 (включно), причому спочатку за спаданням записано всі, що діляться на 5, а потім решту за спаданням.
%enditem




































































































%end
\newpage
\end{document}


